%!TEX root = ../main.tex

\chapter{Retrieval-Augmented Generation}\label{ch:rag}

This chapter describes the \gls{rag} pipeline developed in this work and presents
the experimental results. Section~\ref{sec:rag:arch} describes the architecture
and indexing strategy. Sections~\ref{sec:rag:retrieval}--\ref{sec:rag:reranker}
describe the retrieval and reranking components. Section~\ref{sec:rag:generation}
covers the generation prompt design. Sections~\ref{sec:rag:fullsplit}
and~\ref{sec:rag:dev500} present results on the full validation split and the
dev-500 subset respectively. Section~\ref{sec:rag:analysis} analyses the findings.


% ===========================================================================
\section{Pipeline Architecture}\label{sec:rag:arch}

The \gls{rag} pipeline consists of three sequential stages:

\begin{enumerate}
    \item \textbf{Query preparation.} The raw question is optionally rewritten
          into a self-contained form using the dialogue history
          (Section~\ref{sec:rag:rewrite}).
    \item \textbf{Retrieval.} The (rewritten) query is used to retrieve the
          top-$K$ document sections from the indexed corpus
          (Section~\ref{sec:rag:retrieval}). Optionally, the top-$K'$ retrieved
          candidates are reranked by a cross-encoder and truncated to the top-$K$
          (Section~\ref{sec:rag:reranker}).
    \item \textbf{Generation.} The retrieved sections are formatted as a context
          block and appended to the generation prompt together with the original
          question; the \gls{llm} produces the final answer
          (Section~\ref{sec:rag:generation}).
\end{enumerate}

\subsection{Section-Level Indexing}

Documents are indexed at the \emph{section level}: each section is defined as a
heading together with all content spans that follow it until the next heading,
regardless of heading level. The section-level granularity was chosen to balance
retrieval specificity (coarser than paragraph-level) with context completeness
(finer than document-level).

Each section is stored as a dictionary with fields \texttt{doc\_id},
\texttt{id\_sp} (the span identifier used in grounding annotations),
\texttt{title}, \texttt{section\_title}, \texttt{text}, and
\texttt{span\_ids} (the set of all span IDs covered by the section). The
\texttt{span\_ids} field enables Recall@$K$ to be computed correctly even when a
section covers multiple annotated spans.

The full corpus of 4{,}114 sections is indexed at the start of each experiment;
the BM25 index is built in memory using \texttt{rank\_bm25}, the dense index is
embedded using \texttt{BAAI/bge-small-en-v1.5} and the result matrix is cached
on disk so that subsequent runs skip re-encoding.


% ===========================================================================
\section{Retrieval Strategies}\label{sec:rag:retrieval}

Four retrieval strategies are evaluated. All strategies return the top-$K = 5$
sections for passage to the generator, with the exception of the reranker pipeline
which first retrieves $K' = 50$ candidates.

\subsection{BM25}

The BM25 retriever (\texttt{BM25Okapi}, $k_1 = 1.2$, $b = 0.75$) is built over
the concatenation of section heading and body text, tokenised by whitespace after
lowercasing. All sections are ranked by BM25 score and the top-$K$ are returned.
Because BM25 scores are zero for sections that share no term with the query,
the top-1 score is a potential abstention signal for certain unanswerable
questions~\cite{robertson2009bm25}.

\subsection{Dense Retrieval}

The dense retriever uses \texttt{BAAI/bge-small-en-v1.5} to produce 384-dimensional
embeddings for all sections and for each query. Retrieval is by cosine similarity.
The model was pre-trained on a large collection of sentence pairs and fine-tuned for
general-purpose retrieval, making it competitive with larger models at lower
computational cost.

\subsection{Hybrid Retrieval with RRF}

The hybrid retriever fuses the ranked lists of BM25 and dense retrieval using
\gls{rrf} with $k_{\mathrm{rrf}} = 60$. It returns the top-$K$ sections
by fused score. The RRF formula down-weights sections that rank poorly in either
list while rewarding those that rank highly in both, producing a fused ranking
that is more robust than either individual retriever.

\subsection{Query Rewriting for Conversational Input}\label{sec:rag:rewrite}

Follow-up questions in a multi-turn dialogue are frequently under-specified (e.g.\
``Can I appeal that?'' after a prior turn about a traffic violation decision).
Submitting such questions as-is to a retrieval system typically yields poor results
because the query text contains little discriminating content.

To address this, we implement a \emph{query rewriter}: an additional \gls{llm} call
that reformulates the current question into a self-contained query given the dialogue
history. The rewrite prompt uses a sliding window of the last six dialogue turns
(three user--agent exchanges) and instructs the model to produce a single, complete question that includes
all necessary context. An example rewrite:

\begin{quote}
\small
\textit{History:} User: ``I got a traffic violation.'' Agent: ``I can help with
that. What kind of violation was it?'' User: ``It was a speeding ticket.''

\textit{Original question:} ``Can I appeal that?''

\textit{Rewritten:} ``How can I appeal a speeding ticket traffic violation?''
\end{quote}

Only questions in turns with a non-empty dialogue history are submitted to the
rewriter (434 of 500 dev-500 examples, of which 431 were actually changed);
first-turn questions are passed to retrieval unchanged.


% ===========================================================================
\section{Cross-Encoder Reranking}\label{sec:rag:reranker}

As a second-stage refinement, we evaluate \texttt{BAAI/bge-reranker-base}, a
cross-encoder that scores (query, section) pairs by computing cross-attention
between query tokens and section tokens. This produces more accurate relevance
judgements than bi-encoder cosine similarity at the cost of $O(K')$ forward passes.

The reranking pipeline:
\begin{enumerate}
    \item Retrieve the top $K' = 50$ candidates using the Hybrid+Rewrite retriever.
    \item Score all 50 (rewritten query, section text) pairs with the cross-encoder.
    \item Select the top $K = 5$ sections by cross-encoder score.
\end{enumerate}

Cross-encoder scores for each (query, text) pair are cached on disk
(keyed by SHA-256 of the concatenation), so subsequent runs with the same rewritten
queries are fully cached.

\textbf{Truncation.} The cross-encoder truncates inputs exceeding 512 tokens.
In the first run on dev-500, 1{,}924 of 25{,}000 (query, section) pairs exceeded
this limit; on subsequent cached runs, truncation counts are not recomputed.


% ===========================================================================
\section{Generation Prompt}\label{sec:rag:generation}

The default generation prompt (\emph{v1}) formats the top-$K$ retrieved sections
as a numbered list of context blocks, each headed by document title and section
heading, and instructs the model to answer in one to two plain sentences using only
the provided context:

\begin{quote}
\small\ttfamily
You are a government services agent. Answer the customer's question in 1-2 plain
sentences using only the context below. No bullet points, no bold text, no preamble.
If the context does not contain relevant information, reply exactly:
``Unfortunately, no relevant information is found.''

Examples:\\
Customer: Do I need insurance to register my vehicle?\\
Agent: Yes, you must have liability insurance to register a vehicle in New York.\\
...

Context:\\
\{context\}

Customer: \{question\}\\
Agent:
\end{quote}

\subsection{Agent-Style Prompt (v2)}\label{sec:rag:prompt}

A second prompt variant (\emph{v2}) was designed to reinforce three stylistic
properties that the v1 prompt only weakly enforces: explicit second-person address,
prohibition of markdown formatting, and prohibition of preamble phrases such as
``Based on the context\ldots''. Three few-shot examples were drawn from the
MultiDoc2Dial training split (seed~42, one from each of the DMV, SSA, and VA
domains) to demonstrate the target style. The abstention phrase is unchanged.

The v2 prompt preamble reads:

\begin{quote}
\small\ttfamily
You are a government services agent. Answer the customer's question in 1--2 short,
plain sentences using only the context below. Address the customer directly using
second person (you/your). Do not use bullet points, bold text, or numbered lists.
Do not begin your reply with phrases like ``Based on the context'' or ``According
to the document''.
\end{quote}

The paired bootstrap test between v2 and v1 is reported in
Section~\ref{sec:rag:dev500}.


% ===========================================================================
\section{Oracle Upper Bound}\label{sec:rag:oracle}

To disentangle retrieval failures from generation failures, we establish an
\emph{oracle} baseline in which the model is given the gold grounding sections
directly, bypassing retrieval entirely. Recall@5 is set to 1.0 for all answerable
examples by construction; unanswerable examples receive an empty context.

The oracle \fone{} establishes the \emph{generation ceiling}: the highest \fone{}
achievable by any retrieval system on this data and model combination, assuming
perfect retrieval. Any gap between the best \gls{rag} system and the oracle
represents an upper bound on the gain achievable through retrieval improvements.


% ===========================================================================
\section{Results: Full Validation Split}\label{sec:rag:fullsplit}

Table~\ref{tab:rag:fullsplit} reports generation quality on the full MultiDoc2Dial
validation split ($n = 4{,}427$) for the five original pipelines.
Bootstrap 95\% confidence intervals use 1{,}000 resamples with percentile method
(seed~42).

\begin{table}[ht]
    \centering
    \caption{Generation quality on the full validation split ($n = 4{,}427$).
             95\% bootstrap CIs in brackets.}
    \label{tab:rag:fullsplit}
    \begin{tabular}{lrr}
        \toprule
        System & EM\% [CI] & \fone{}\% [CI] \\
        \midrule
        LLM-only           & 0.77 [0.54, 1.04] & 15.88 [15.54, 16.29] \\
        RAG BM25           & 2.67 [2.19, 3.14] & 24.11 [23.46, 24.77] \\
        RAG Dense          & 3.52 [2.98, 4.07] & 24.58 [23.89, 25.28] \\
        RAG Hybrid         & 2.67 [2.22, 3.19] & 24.75 [24.09, 25.43] \\
        RAG Hybrid+Rewrite & 1.90 [1.52, 2.31] & 26.78 [26.19, 27.46] \\
        \bottomrule
    \end{tabular}
\end{table}

A paired bootstrap test on the full split (same procedure as
Section~\ref{sec:rag:dev500}, restricted to the $4{,}419$ examples for which both
systems returned an answer) confirms the gain from query rewriting: Hybrid+Rewrite
outperforms Hybrid by $\Delta\fone = {+}2.03$~pp (95\% CI $[1.43, 2.57]$,
$p < 0.001$), and outperforms the \gls{llm}-only baseline by
$\Delta\fone = {+}10.89$~pp (95\% CI $[10.28, 11.50]$, $p < 0.001$).

Table~\ref{tab:rag:fullsplit_ret} reports retrieval quality. Domain Recall@5 measures
the fraction of the top-5 sections drawn from the correct domain.

\begin{table}[ht]
    \centering
    \caption{Retrieval quality on the full validation split.
             95\% bootstrap CIs in brackets.}
    \label{tab:rag:fullsplit_ret}
    \begin{tabular}{lrrr}
        \toprule
        System & \ratk{1}\% [CI] & \ratk{5}\% [CI] & Domain \ratk{5}\% [CI] \\
        \midrule
        RAG BM25           & 21.03 [19.83, 22.18] & 40.61 [39.32, 42.01] & 65.95 [64.86, 67.04] \\
        RAG Dense          & 21.51 [20.24, 22.74] & 40.96 [39.48, 42.50] & 70.46 [69.34, 71.65] \\
        RAG Hybrid         & 23.85 [22.56, 25.07] & 44.70 [43.21, 46.19] & 69.39 [68.35, 70.43] \\
        RAG Hybrid+Rewrite & 30.64 [29.35, 31.98] & 59.45 [58.04, 60.87] & 92.46 [91.85, 93.04] \\
        \bottomrule
    \end{tabular}
\end{table}

Table~\ref{tab:rag:docswitch} shows retrieval quality broken down by dialogue turn
type (document-switch vs.\ same-document).

\begin{table}[ht]
    \centering
    \caption{Recall@5 by turn type on the full validation split.}
    \label{tab:rag:docswitch}
    \begin{tabular}{lrr}
        \toprule
        System & Doc-switch \ratk{5}\% & Same-doc \ratk{5}\% \\
        \midrule
        RAG BM25           & 50.28 & 38.12 \\
        RAG Dense          & 55.00 & 37.35 \\
        RAG Hybrid         & 56.59 & 41.65 \\
        RAG Hybrid+Rewrite & 52.71 & 61.18 \\
        \bottomrule
    \end{tabular}
\end{table}


% ===========================================================================
\section{Results: Dev-500 Diagnostic Subset}\label{sec:rag:dev500}

Table~\ref{tab:rag:dev500} reports the full diagnostic evaluation on dev-500,
including the oracle ceiling and the v2 prompt and reranker ablations.
All systems use $K = 5$ retrieved sections.

\begin{table}[htbp]
    \centering
    \small
    \caption{Diagnostic evaluation on dev-500 ($n = 500$, 95\% bootstrap CI).
             AR = abstain recall, FA = false-abstention rate,
             AP = abstain precision.}
    \label{tab:rag:dev500}
    \begin{tabular}{lrrrrr}
        \toprule
        System & \fone{}\% [CI] & \ratk{5}\% & AR\% & FA\% & AP\% \\
        \midrule
        LLM-only           & 16.3 [15.0, 17.6] &  ---  & 19.4 &  7.7 & 14.3 \\
        Oracle             & 36.1 [34.0, 38.4] & 100.0 & 25.8 &  4.1 & 29.6 \\
        Oracle (v2 prompt) & 36.3 [34.4, 38.5] & 100.0 & 16.1 &  1.9 & 35.7 \\
        \midrule
        RAG BM25           & 24.5 [22.4, 26.4] &  39.1 & 41.9 & 15.4 & 15.3 \\
        RAG Dense          & 25.7 [23.5, 28.0] &  38.4 & 83.9 & 22.1 & 20.2 \\
        RAG Hybrid         & 26.1 [24.1, 28.1] &  43.2 & 61.3 & 16.0 & 20.2 \\
        RAG Hybrid+Rewrite & 28.0 [26.2, 30.0] &  61.2 & 35.5 &  7.0 & 25.0 \\
        \quad (v2 prompt)  & 27.8 [25.9, 29.8] &  61.2 & 29.0 &  4.9 & 28.1 \\
        \midrule
        Hybrid+Rewrite+Rerank           & 28.2 [26.3, 30.0] &  67.4 & 32.3 &  6.8 & 23.8 \\
        \quad (v2 prompt)  & 28.0 [26.2, 29.9] &  67.4 & 25.8 &  5.1 & 25.0 \\
        \bottomrule
    \end{tabular}
\end{table}

\subsection{Paired Bootstrap Tests}

Table~\ref{tab:rag:paired} reports paired bootstrap tests on token \fone{} between
selected system pairs. The test statistic is the mean paired difference; the
$p$-value is the fraction of bootstrap resamples in which the difference is
$\leq 0$ (i.e.\ one-sided test: does system A outperform system B?).

\begin{table}[ht]
    \centering
    \caption{Paired bootstrap tests on token \fone{} (dev-500, $n = 500$, $B = 1{,}000$).}
    \label{tab:rag:paired}
    \begin{tabular}{llrrr}
        \toprule
        System A & System B & $\Delta$\fone{}\% & 95\% CI & $p$-value \\
        \midrule
        Hybrid+Rewrite      & Hybrid            & $+1.94$ & $[0.07, 3.48]$  & 0.017$^*$ \\
        Oracle (v2)         & Oracle            & $+0.25$ & $[-0.87, 1.32]$ & 0.319 \\
        Hybrid+Rewrite (v2) & Hybrid+Rewrite    & $-0.21$ & $[-1.68, 1.35]$ & 0.607 \\
        Rerank              & Hybrid+Rewrite    & $+0.16$ & $[-1.27, 1.60]$ & 0.414 \\
        Rerank (v2)         & Hybrid+Rewrite (v2) & $+0.22$ & $[-1.13, 1.56]$ & 0.394 \\
        \bottomrule
        \multicolumn{5}{l}{\small $^*$ statistically significant at $\alpha = 0.05$}
    \end{tabular}
\end{table}


% ===========================================================================
\section{Analysis}\label{sec:rag:analysis}

\subsection{The Retrieval--Generation Gap}

The oracle ceiling of $\fone = 36.1\%$ on dev-500 reveals that even with perfect
retrieval the model produces answers that overlap only about one-third of the
reference tokens. This reflects the nature of the task: reference answers in
\gls{mdd} are verbatim extracts from the source document, while the model is
instructed to produce concise, agent-style paraphrases. Token \fone{} undervalues
semantically correct paraphrases, so the true generation quality may be higher than
it appears.

The best \gls{rag} system (Hybrid+Rewrite) achieves $\fone = 28.0\%$, an 8.1~pp
gap below the oracle ceiling. On dev-500, $286$ of $469$ answerable examples have
the gold section in the top-5 (retrieval hit), and F1-hit for those examples is
33.3\%. For the $183$ examples (39\% of answerable ones) where retrieval misses,
$\fone = 18.3\%$: the model is limited to whatever partial context it retrieves.
This \emph{hit--miss gap} is large for every retrieval strategy: 21.8~pp for BM25,
22.3~pp for Dense, 21.5~pp for Hybrid, and 15.0~pp for Hybrid+Rewrite.

These figures allow the oracle gap to be decomposed. If every retrieval miss were
converted into a hit at the observed F1-hit, \fone{} would rise by approximately
$183 \times (33.3 - 18.3) / 500 \approx 5.5$~pp, i.e.\ roughly two thirds of the
8.1~pp gap. The remainder reflects the lower F1-hit of the Hybrid+Rewrite system
(33.3\%) relative to the oracle (36.7\%), plausibly caused by distractor
sections in the top-5 context, which the oracle does not contain. Retrieval
coverage therefore remains the primary lever for improving answer quality.

\subsection{Effect of Query Rewriting}

Query rewriting produces the largest single improvement across all evaluated
interventions: $+1.94$~pp \fone{} ($p = 0.017$) on dev-500 and $+2.03$~pp
($p < 0.001$) on the full validation split. More strikingly, \ratk{5} rises from 43.2\% (Hybrid, no
rewrite) to 61.2\% ($+18$~pp), and domain recall rises to 92.5\% --- the rewriter
effectively anchors retrieval to the correct domain in almost all cases.

The improvement is concentrated on same-document follow-up questions, where the
decontextualised rewrite provides the exact vocabulary of the target document
section. Document-switch turns, by contrast, benefit less from rewriting because
the rewriter encodes the vocabulary of the \emph{previous} document, pulling
retrieval toward it and away from the new one. Hybrid retrieval (no rewrite)
achieves better doc-switch \ratk{5} (56.6\%) than Hybrid+Rewrite (52.7\%).

\subsection{Effect of Cross-Encoder Reranking}

The cross-encoder reranker improves \ratk{5} from 61.2\% to 67.4\% ($+6.2$~pp) and
increases the retrieval-hit count from 286 to 316 examples ($+30$). Despite this,
\fone{} improves by only $+0.16$~pp ($p = 0.414$, not significant).

This small aggregate change is \emph{not} explained by a generation ceiling. Given
the hit--miss gap of $\sim$15~pp, a net gain of 30 hits would be expected to raise
\fone{} by roughly $30 \times 15 / 500 \approx 0.9$~pp. Table~\ref{tab:rag:rerank}
decomposes the observed change by how the reranker altered each example's
retrieval outcome.

\begin{table}[ht]
    \centering
    \caption{Per-example decomposition of the \fone{} change from Hybrid+Rewrite
             to Hybrid+Rewrite+Rerank (dev-500, v1 prompt). ``Contribution'' is
             the group's share of the change in the 500-example mean.}
    \label{tab:rag:rerank}
    \small
    \begin{tabular}{lrrrr}
        \toprule
        Retrieval outcome (Rewrite $\to$ Rerank) & $n$ & \fone{} before & \fone{} after & Contribution (pp) \\
        \midrule
        miss $\to$ hit   &  53 & 21.9 & 33.2 & $+1.20$ \\
        hit $\to$ miss   &  23 & 30.4 & 21.3 & $-0.42$ \\
        hit $\to$ hit    & 263 & 33.6 & 33.5 & $-0.06$ \\
        miss $\to$ miss  & 130 & 16.9 & 16.0 & $-0.21$ \\
        unanswerable     &  31 &  --- &  --- & $-0.35$ \\
        \midrule
        Total            & 500 &      &      & $+0.16$ \\
        \bottomrule
    \end{tabular}
\end{table}

When the reranker newly retrieves the gold section (53 examples), \fone{} rises by
11.3~pp to 33.2\%, matching the F1-hit of examples that were already hits. The
generator therefore exploits the additional evidence fully. The aggregate gain is
small because the reranker also \emph{displaces} the gold section for 23 examples
that were previously hits (\fone{} $-9.1$~pp), slightly degrades the
miss$\to$miss group, and abstains less often on unanswerable questions (abstain
recall 35.5\% $\to$ 32.3\%). The net $+30$ hits is the difference of 53 gains and
23 losses, and the resulting effect is too small to be detected at $n = 500$.
Reranking thus remains a promising direction, provided the loss of previously
correct sections can be reduced (e.g.\ by fusing the cross-encoder score with the
first-stage RRF score rather than replacing it).

\subsection{Effect of the Agent-Style Prompt}

The v2 prompt produces no statistically significant change in \fone{} on either the
oracle or the Hybrid+Rewrite system. However, it consistently reduces the
false-abstention rate: from 4.1\% to 1.9\% on Oracle, and from 7.0\% to 4.9\% on
Hybrid+Rewrite. The model becomes more conservative --- it abstains less frequently
overall but more accurately. The tradeoff is a reduction in abstain recall (Oracle:
25.8\% $\to$ 16.1\%; Hybrid+Rewrite: 35.5\% $\to$ 29.0\%).

\subsection{Abstention Calibration}

Among the five original retrieval strategies, Dense achieves the highest abstain
recall (83.9\%) but at the cost of a 22.1\% false-abstention rate --- it abstains
on far too many answerable questions. Its abstain precision (20.2\%) is no higher
than Hybrid's, which indicates that Dense does not identify unanswerable questions
more accurately; it simply abstains more often overall. A plausible explanation is
that dense retrieval returns topically related but non-answering sections more
often than lexical retrieval does (its \ratk{5} is the lowest of the four
strategies), so the generator more frequently judges the context insufficient and
falls back to the abstention phrase, for answerable and unanswerable questions
alike.

Hybrid+Rewrite achieves the best calibration among the five systems: 35.5\% abstain
recall, 7.0\% false abstentions, and 25.0\% abstain precision. The v2 prompt shifts
this to 29.0\% / 4.9\% / 28.1\%, trading some recall for a lower false-abstention
rate and higher precision.

Improving abstain recall without increasing false abstentions remains an open
problem for this pipeline. Approaches such as a retrieval sufficiency check (an
additional \gls{llm} call that judges whether the retrieved context answers the
question) or training a dedicated classification head are planned for future work.
